%! Polynomial Equations, Inequalities, and Modeling — Unit 2, Lesson 2.7 — Student Packet Cover Sheet
\documentclass[10pt]{article}
\usepackage{precalculus-article}
\usepackage{precalculus-boxes}
\usepackage{ltablex}
\keepXColumns

\begin{document}

% Full-bleed plum banner
\begin{tikzpicture}[remember picture, overlay]
  \fill[plum] (current page.north west)
    rectangle ([yshift=-0.9in]current page.north east);
\end{tikzpicture}
\vspace{-0.8in}

\begin{center}
  {\color{white}\textbf{\LARGE Precalculus}}\\[4pt]
  {\color{lilacmid}\large\textbf{Unit 2: Polynomial Functions}}\\[6pt]
  {\color{white}\textbf{\large Lesson 2.7 \quad Polynomial Equations, Inequalities, and Modeling}}
\end{center}

\vspace{0.22in}
\namedateperiod
\vspace{0.18in}

\begin{learningtargetbox}
\small
\begin{enumerate}[leftmargin=1.5em, itemsep=5pt]
  \item I can solve a polynomial equation by collecting every term on one side
        and applying the zero-product property, and I can explain why dividing
        both sides by $x$ throws a solution away.
  \item I can solve a polynomial inequality by marking its real zeros on a
        number line, testing one value in each interval, and writing the
        solution in interval notation --- using multiplicity to predict where
        the sign changes and where it does not.
  \item I can construct a polynomial model from a described scenario and state
        the domain restriction the context forces, and I can say when a
        regression is the right tool instead.
  \item I can use a model to answer a question about the scenario --- including
        an ``at least'' question that becomes an inequality --- and report the
        answer in a sentence with units.
\end{enumerate}
\end{learningtargetbox}

\vspace{0.14in}

\begin{tocbox}
\small
\renewcommand{\arraystretch}{1.5}
\begin{tabularx}{\linewidth}{@{} c l X r @{}}
\rowcolor{lilac}
\textbf{\#} & \textbf{Component} & \textbf{Description} & \textbf{Score} \\
\midrule
1 & Warm-Up        & Spiral review of prerequisite skills               & \blank{1.2cm} \\
2 & Guided Notes   & Equations, sign charts, and building a model       & \blank{1.2cm} \\
3 & Group Activity & The Packaging Desk: one sheet, three tiers         & \blank{1.2cm} \\
4 & Exit Ticket    & Independent check on equations, signs, and domain  & \blank{1.2cm} \\
5 & Homework       & Practice and extension                             & \blank{1.2cm} \\
\midrule
  & \multicolumn{2}{r}{\textbf{Total}} & \blank{1.2cm} \\
\end{tabularx}
\end{tocbox}

\vspace{0.14in}

\begin{remindbox}
\small
Everything this unit built comes due today. A solution of $p(x) = 0$, a real
zero of $p$, and an $x$-intercept of the graph have been the same object under
three names since Lesson 2.3 --- and now those zeros take on a new job. They
cut the number line into pieces, and on each piece the sign of $p$ cannot
change. That is what turns a question like ``which cuts give a box holding at
least $1000$ cubic centimeters?'' into something you can actually answer.
\end{remindbox}

\end{document}
